\thispagestyle{empty}%
%
\begin{center}
    \LARGE\textbf{\scshape{\OlympString{tournamelowercase}}}
\end{center}
%
\vspace{2pt}%
%
\begin{center}
    \large{\OlympString{touronedate}}
\end{center}
\vspace{2pt}
\begin{center}
    \large\textbf{Сначала, пожалуйста, прочитайте следующее:}
\end{center}
\vspace{4pt}

\begin{itemize}[leftmargin=*]
    \item Данный тур состоит из \OlympString{problemcount_instructions}.
    \item Необходимые для Вашего решения математические подсказки и список физических констант предоставлены внутри соответствующих задач.
    \item Всего страниц в данном комплекте страниц на Вашем языке ---  \pageref*{LastPage}, \textit{исключая} титульный лист и страницу с инструкциями, которую Вы читаете сейчас.
    \item Для тура Вам потребуются канцелярские принадлежности и инженерный калькулятор.
    \item Вам будут предоставлены чистые листы бумаги для черновых записей и \textbf{\textit{Поля для заполнения решений}} (чистовики). При нехватке бумаг поднимите руку и сообщите подошедшему к Вам наблюдателю --- Вам выдадут дополнительные листы.
    \item Ваша работа будет оцениваться по \textbf{\textit{чистовикам}}; записи на черновиках оцениваться не будут. В основном Вы должны использовать уравнения, числа, буквенные обозначения, рисунки и графики; словесные описания не оцениваются, если в условии об этом не спрашивается.
    \item При оформлении понятно указывайте номер задачи, которую Вы решаете. Решение каждой задачи следует начинать с новой страницы \textbf{\textit{чистовиков}}.
    \item Используйте только лицевую сторону \textbf{\textit{чистовиков}} и постарайтесь не выходить за пределы рамок. В верхней части \textbf{\textit{чистовиков}} Вы должны указать \textit{порядковый номер страницы}. Если Вы не хотите, чтобы определённые листы чистовиков учитывались в оценивании, то перечеркните их большим крестом на весь лист и не включайте в Ваш подсчёт полного количества листов.
    \item \textbf{Не указывайте информацию о себе в чистовиках (имя, фамилия, школа, и т.д.).} Ваши работы будут зашифрованы организаторами по окончании тура.
    \item Дополнительные советы:
    \begin{itemize}[label=$\circ$]
        \item Убедитесь, что Вы прочитали условия до конца. Иногда можно пропустить некоторые пункты задач и успешно решить последующие.
        \item Как можно подробнее выписывайте все нужные уравнения, даже если Вы не смогли прийти к конечному ответу.
        \item Не пренебрегайте возможностью перепроверять свои решения, а также проверять ответ и промежуточные формулы на размерность --- небольшая ошибка может Вам стоить большого количества баллов!
        \item Оформляйте попытки решения сразу на \textbf{\textit{чистовики}}, а черновики используйте только для промежуточных расчётов, решения громоздких систем уравнений, и алгебраических упрощений. Вам нужно постараться изложить все Ваши решения в чистовик, не тратя лишнего времени на переписывание из черновика.
        \item Используйте отведенное время полностью, даже если Вам кажется, что никаких идей больше нет. Очень часто правильный метод решения приходит после долгого обдумывания. Также важно перечитывать условия после написания решения, чтобы убедиться, что найденная Вами величина именно та, которую требовалось найти.
    \end{itemize}
\end{itemize}
%
\vspace{4pt}%
%
\begin{center}
    \textsc{Желаем Вам удачи!}
\end{center}%
%
\begin{center}
    \textit{\OlympString{duration}}
\end{center}

\newpage